\addcontentsline{toc}{subsubsection}{Solution 59. The Hidden Product Rule}

\begin{dlsummary}[title={\bfseries Solution 59. The Hidden Product Rule},breakable]

\begin{proof}
The unique solution is
\[
\boxed{f(x)=x\qquad(x\in\mathbb R).}
\]

\paragraph{1. Remove the expected solution.}

Define
\[
g(x)=f(x)-x.
\]
The first two inequalities become
\[
g(x+y)\ge g(x)+g(y), \tag{1}
\]
\[
g(xy)\ge xg(y)+yg(x). \tag{2}
\]

We first show that both are actually equalities.

\paragraph{2. Establish the normalisation and oddness.}

Putting \(y=0\) in (2),
\[
g(0)\ge xg(0)\qquad(x\in\mathbb R),
\]
which forces \(g(0)=0\).

Putting \(x=1\) in (2),
\[
0\ge yg(1)\qquad(y\in\mathbb R),
\]
so \(g(1)=0\).

Taking \(x=y=-1\) in (2) gives
\[
0\ge-2g(-1),
\]
so \(g(-1)\ge0\). But (1) at \((1,-1)\) gives \(g(-1)\le0\). Hence \(g(-1)=0\).

Now (2) at \((-1,x)\) gives
\[
g(-x)\ge-g(x),
\]
while (1) at \((x,-x)\) gives the reverse inequality. Therefore
\[
\boxed{g(-x)=-g(x).} \tag{3}
\]

\paragraph{3. Force additivity and the product identity.}

Apply (1) to \(-x,-y\), then use oddness:
\[
g(x+y)\le g(x)+g(y).
\]
Thus
\[
\boxed{g(x+y)=g(x)+g(y).} \tag{4}
\]

Similarly, applying (2) to \(-x,y\) reverses its direction. Hence
\[
\boxed{g(xy)=xg(y)+yg(x).} \tag{5}
\]

\paragraph{4. Interpret the composition inequality.}

Since \(f(t)=t+g(t)\), additivity gives
\[
f(f(t))
=f(t)+g(f(t))
=t+2g(t)+g(g(t)).
\]
Consequently, the third original inequality becomes
\[
\boxed{g(g(t))\ge0\qquad(t\ge0).} \tag{6}
\]

Let \(h=g\circ g\). It is additive, and (6) makes it nondecreasing: if \(b>a\), then
\[
h(b)-h(a)=h(b-a)\ge0.
\]
Moreover,
\[
h(1)=g(g(1))=0.
\]
An additive, nondecreasing function with value \(0\) at \(1\) is identically zero: it vanishes on the rationals, and rational bounds squeeze its value at every real input. Thus
\[
\boxed{g(g(x))=0\qquad(x\in\mathbb R).} \tag{7}
\]

\paragraph{5. Apply the product identity twice.}

From (5),
\[
g(x^2)=2xg(x).
\]
Apply \(g\) again. By additivity and (5),
\[
\begin{aligned}
g(g(x^2))
&=2g(xg(x))\\
&=2xg(g(x))+2g(x)^2.
\end{aligned}
\]
Both composite terms vanish by (7), leaving
\[
0=2g(x)^2.
\]
Hence \(g(x)=0\) everywhere, so \(f(x)=x\).

The identity function satisfies all three inequalities with equality.
\end{proof}
\end{dlsummary}

\begin{takeaways}
\begin{itemize}
\item \textbf{Subtract the expected model:} This exposes the algebraic structure hidden in the inequalities.
\item \textbf{Negative inputs reverse information:} Oddness turns superadditivity and the product inequality into exact identities.
\item \textbf{Regularise the composite first:} Positivity on nonnegative inputs makes \(g\circ g\) monotone, even though no regularity of \(g\) is assumed.
\item \textbf{Apply a structural identity twice:} The final square \(g(x)^2\) eliminates every remaining possibility.
\end{itemize}
\end{takeaways}
